\textbf{Automatic NLG evaluation}
In the landscape of NLG, evaluating the quality of generated text represents a particularly arduous task. For a significant period, evaluation was primarily dependent on human annotations, a process that is labor-intensive and limited by scalability issues. Automatic NLG evaluation attempts to address these challenges by leveraging computational models to assess the quality of a generated text. Previous work lies on the following categories: (1) \textit{n-gram-based metrics}: ROUGE~\citep{lin2004rouge} is a set of metrics that compute the amount of overlap between n-grams in the machine-generated summaries and the reference summaries. BLEU~\citep{papineni2002bleu} compare the generated text with reference translations, based on the co-occurrence of n-grams in both texts. In spite of being easily and widely used, the above method is incapable of capturing syntactic and semantic similarity~\citep{stent2005evaluating}. (2) \textit{embedding-based metrics}: Word embeddings are vector representations of words that capture their semantic properties, such that words with similar meanings have similar embeddings. A bunch of work leverages word embeddings to evaluate the semantic similarity between two pieces of text. BERTScore~\citep{zhang2019bertscore} use contextualized word embeddings from transformer models like BERT~\citep{devlin2018bert}, BLEURT~\citep{sellam2020bleurt} utilize supervised training data to enhance the performance. MoverScore~\citep{zhao2019moverscore} combine contextualized word embeddings with Earth Mover’s Distance~\citep{rubner2000earth}. (3) \textit{LLM-based metrics}: Amidst the flourishing advancement of LLM which embodies a wealth of information derived from extensive training data, using LLM as an evaluator has experienced notable progress. GPTScore~\citep{fu2023gptscore} utilize conditional probability to assign the text a score representing its quality.~\cite{wang2023chatgpt} explore the potential of utilizing ChatGPT as an NLG evaluator by prompting it to score a text directly.~\cite{wang2023pandalm} curate a reliable dataset containing pairwise comparison and evaluation explanation which can be used to train a foundation model making it a better evaluator.~\cite{bai2023benchmarking} propose decentralized 
evaluation to provide fairer evaluation results. G-EVAL~\citep{liu2023gpteval} propose probability-weighted techniques to calibrate the score given by a single LLM. 

\textbf{Communicative Agents} Most recently, significant attention has been dedicated to the development of communicative agents. These agents, often acted by LLMs like ChatGPT or GPT-4, are designed to interact and communicate effectively with other agents or human users using natural language. The primary goal is to facilitate more productive and efficient interaction and collaboration as different agents can autonomously communicate and negotiate to tackle a more complex task collectively. Several studies have explored various aspects of communicative agents.~\cite{li2023camel} propose a cooperative agent framework dubbed as \textit{role-playing} enabling agents to autonomously cooperate to solve complex tasks.~\cite{park2023generative} create a sandbox environment consisting of 25 individual virtual entities endowed with a character description and memory system. Every intelligent agent is capable of autonomously interacting with other agents and the environment simulating reliable human behavior.~\cite{qian2023communicative} establish a chat-based software development framework that can complete a software design and produce executable software at a reduced cost compared to recruiting human programmers.~\cite{liu2023training} utilize a sandbox environment to curate reliable datasets in better alignment with human preference and train a socially-aligned LLM.~\cite{liang2023encouraging} and~\cite{du2023improving} also make use of the multi-agent debate framework in other scenarios such as translation and arithmetic problems resulting in better results.~\cite{wang2023unleashing} propose an alternative method called self-collaboration to enable the communication of agents by utilizing a single LLM prompted by multi-persona descriptions.~\cite{mandi2023roco} propose a novel framework designed for the collaboration of multiple robots, utilizing multiple LLMs to enhance coordination and strategic planning among the robots. Concurrent with our work,~\cite{li2023prd} propose Peer Rank and Discussion (PRD) which is similar to our approach. However, they probe the different dimensions of evaluation by using different models as agents and did not explore alternative communication strategies.
